\begin{proof}[Proof of \cref{obj:thm:gap-free-positive-measure-modulus}]
\begin{enumerate}
\item Put
\[
s_\star:=\pi_0\sigma_0^2,
\qquad
L_\tau:=\frac{4L\sqrt{d_z}}{\sigma_0}.
\]
For matrix sizes \(r_{\mathrm{row}},r_{\mathrm{col}}\), write
\[
\operatorname{Ent}_{r_{\mathrm{row}},r_{\mathrm{col}}}
:=
r_{\mathrm{col}}r_{\mathrm{row}}\,
\|J_{r_{\mathrm{col}}}\|_{\mathrm{op}},
\]
with \(J_{r_{\mathrm{col}}}\) as in \cref{obj:lem:summary-metric-coordinate-comparison}.  Define
\[
\begin{aligned}
K_{\mathrm{forw}}
&:=
\operatorname{Ent}_{d_x,k}\operatorname{Ent}_{k,k}
\operatorname{Ent}_{k,d_x}\,\sigma_0^{-1}
+\operatorname{Ent}_{d_x,d_x}
+\operatorname{Ent}_{d_x,k}\operatorname{Ent}_{k,d_x},\\
K_{\mathrm{back}}
&:=
\operatorname{Ent}_{d_x,k}\operatorname{Ent}_{k,k}
\operatorname{Ent}_{k,d_x}\,(d_xL)
+\operatorname{Ent}_{d_x,d_x}
+\operatorname{Ent}_{d_x,k}\operatorname{Ent}_{k,d_x},
\end{aligned}
\]
and
\[
\kappa_\star
:=
\operatorname{Ent}_{d_x,d_x}^{\,2}K_{\mathrm{forw}}K_{\mathrm{back}} .
\]
The assumptions \(L\ge1\) and \(\sigma_0>0\), together with nonnegativity of the entry-norm constants, give
\[
0\le \kappa_\star .
\]
Set
\[
C_{\mathrm{base}}
:=
\kappa_\star L_\tau
+
L\,d_x^2\kappa_\star^2
\{3s_\star^{-2}L+s_\star^{-1}\},
\qquad
C_{\mathrm{mod}}:=C_{\mathrm{base}}+1 .
\]
Since \(s_\star>0\), \(L_\tau\ge0\), and \(\kappa_\star\ge0\), we have
\[
0\le C_{\mathrm{base}},
\qquad
0<C_{\mathrm{mod}} .
\]
% lean: modelRealDiagonalization_conditionBound_nonneg

\item Fix \(P\in\mathcal M\).  Let \(B(P)\) be the target-proxy feature matrix of \cref{obj:def:target-feature}.  The proxy-rank margin gives
\[
\sigma_0\le \sigma_k(B(P)).
\]
Choose a thin singular-value factorization
\[
B(P)=V_P\Upsilon_P,
\qquad
V_P^\top V_P=I_k,
\qquad
\Upsilon_P^{-1}\Upsilon_P=\Upsilon_P\Upsilon_P^{-1}=I_k .
\]
The singular-value lower bound implies
\[
|(\Upsilon_P^{-1})_{ab}|\le \sigma_0^{-1}
\qquad(1\le a,b\le k).
\]
For every coordinate \(i\) and latent class \(u\), latent-arm positivity gives
\[
0<P(U=u,T=0)\le P(U=u).
\]
Under the normalized restriction of \(P\) to \(\{U=u\}\), \cref{obj:lem:proxy-coordinate-envelopes} gives
\[
|X_i|\le L\quad\text{almost surely}.
\]
Therefore
\[
\begin{aligned}
|B(P)_{iu}|
&=\left|E_P[X_i\mid U=u]\right|\\
&=\left|\int X_i\,dP(\,\cdot\mid U=u)\right|\\
&\le \int |X_i|\,dP(\,\cdot\mid U=u)
\le L .
\end{aligned}
\]
Thus \cref{obj:lem:diagonalization-condition-number-bound} applies to \(B(P)=V_P\Upsilon_P\).

Define the ambient contrast
\[
\mathfrak D_P
:=
M_1(P)^\dagger N_1(P)-M_0(P)^\dagger N_0(P).
\]
By \cref{obj:lem:observed-proxy-moment-factorization,obj:lem:observed-outcome-proxy-moment-factorization}, for each \(t\in\{0,1\}\),
\[
M_t(P)=A_t(P)\Gamma_t(P)B(P)^\top,
\qquad
N_t(P)=A_t(P)\Gamma_t(P)
\operatorname{diag}\{\mu_{tu}(P)\}_{u=1}^kB(P)^\top ,
\]
where
\[
\Gamma_t(P)=
\operatorname{diag}\left\{
\frac{P(U=u,T=t)}{P(T=t)}:1\le u\le k
\right\}.
\]
The proxy-rank margin and \cref{obj:lem:latent-arm-weight-invertibility} make
\(A_t(P)\Gamma_t(P)\) injective, and the proxy-rank margin makes \(B(P)\) injective.  Hence
\cref{obj:lem:moore-penrose-product-cancellation} yields
\[
M_t(P)^\dagger N_t(P)
=
(B(P)^\dagger)^\top
\operatorname{diag}\{\mu_{tu}(P)\}_{u=1}^kB(P)^\top .
\]
Subtracting the two arms gives
\[
\mathfrak D_P
=
(B(P)^\dagger)^\top
\operatorname{diag}\{\tau_u(P)\}_{u=1}^kB(P)^\top .
\]

Let
\[
\mathsf{Forw}_P
:=
V_P(\Upsilon_P^{-1})^\top V_P^\top+(I_{d_x}-V_PV_P^\top),
\qquad
\mathsf{Back}_P
:=
V_P\Upsilon_P^\top V_P^\top+(I_{d_x}-V_PV_P^\top).
\]
Then
\[
\mathsf{Forw}_P\mathsf{Back}_P
=
\mathsf{Back}_P\mathsf{Forw}_P
=
I_{d_x}.
\]
Extending the columns of \(V_P\) to an orthonormal basis of \(\mathbb R^{d_x}\), the columns obtained by applying \(\mathsf{Forw}_P\) form a real eigenbasis for \(\mathfrak D_P\).  The signal eigenvalues are \(\tau_u(P)\), and the remaining eigenvalues are \(0\).  Consequently
\begin{equation}
\mathfrak D_P
=
\mathcal S_P\operatorname{diag}(\alpha_{P,1},\ldots,\alpha_{P,d_x})
\mathcal S_P^{-1},
\label{eq:gap-free-model-diagonalization}
\end{equation}
with
\[
\|\mathcal S_P\|_{\mathrm{op}}\|\mathcal S_P^{-1}\|_{\mathrm{op}}
\le \kappa_\star,
\qquad
|\alpha_{P,j}|\le L_\tau .
\]
The eigenvalue bound uses \cref{obj:lem:latent-effect-envelope} on the signal coordinates and \(0\in[-L_\tau,L_\tau]\) on the remaining coordinates.

It remains to record the two anchor identities used by the spectral representation.  For each coordinate \(j\),
\[
\begin{aligned}
m_X(P)_j
&=\int X_j\,dP\\
&=\sum_{u=1}^k\int_{\{U=u\}}X_j\,dP\\
&=\sum_{u=1}^k E_P[X_j\mid U=u]\,P(U=u)\\
&=\sum_{u=1}^k B(P)_{ju}p_u(P),
\end{aligned}
\]
where the latent classes form a finite measurable partition and \(X_j\) is integrable by \cref{obj:lem:proxy-coordinate-envelopes}.  Hence
\[
m_X(P)=B(P)p(P).
\]
Also, since \(X_1=1\) almost surely and \(P(U=u)>0\),
\[
B(P)_{1u}=E_P[X_1\mid U=u]=1
\qquad(1\le u\le k),
\]
so the first canonical basis vector \(e_1\in\mathbb R^{d_x}\) satisfies
\[
B(P)^\top e_1=\mathbf 1_k .
\]
Therefore, for every one-Lipschitz \(\varphi:\mathbb R\to\mathbb R\) satisfying
\(\varphi(0)=0\),
\[
\varphi(\mathfrak D_P)
=
(B(P)^\dagger)^\top
\operatorname{diag}\{\varphi(\tau_u(P))\}_{u=1}^kB(P)^\top
\]
and
\[
\int \varphi\,d\nu_P
=
\sum_{u=1}^kp_u(P)\varphi(\tau_u(P))
=
\left\langle m_X(P),\varphi(\mathfrak D_P)e_1\right\rangle .
\]
% lean: model_realDiagonalization_certificate

\item Fix \(P,Q\in\mathcal M\).  For each arm \(t\), \cref{obj:prop:observed-vmw-margin-inclusion} supplies
\[
s_\star\le\sigma_k(M_t(P)),
\qquad
s_\star\le\sigma_k(M_t(Q)),
\qquad
\|N_t(P)\|_{\mathrm{op}}\le L .
\]
Together with the proxy factorization, these inequalities give
\[
\operatorname{rank}M_t(P)=\operatorname{rank}M_t(Q)=k .
\]
We first record the armwise perturbation estimate.  Let \(\mathsf H,\mathsf H'\in\mathbb R^{d_z\times d_x}\) have the same rank \(k\), and suppose
\[
s_\star\le \sigma_k(\mathsf H),
\qquad
s_\star\le \sigma_k(\mathsf H') .
\]
Then
\[
\|\mathsf H^\dagger\|_{\mathrm{op}}\le s_\star^{-1},
\qquad
\|\mathsf H'{}^\dagger\|_{\mathrm{op}}\le s_\star^{-1},
\]
because the operator norm of the Moore--Penrose inverse is the reciprocal of the smallest nonzero singular value.  The Penrose equations, applied with the common rank condition, give the exact moving-space identity
\[
\begin{aligned}
\mathsf H^\dagger-\mathsf H'{}^\dagger
={}&-\mathsf H^\dagger(\mathsf H-\mathsf H')\mathsf H'{}^\dagger
+\mathsf H^\dagger(\mathsf H^\dagger)^\top
  (\mathsf H-\mathsf H')^\top(I-\mathsf H'\mathsf H'{}^\dagger)\\
&+(I-\mathsf H^\dagger\mathsf H)(\mathsf H-\mathsf H')^\top
  (\mathsf H'{}^\dagger)^\top\mathsf H'{}^\dagger .
\end{aligned}
\]
Here \(I-\mathsf H'\mathsf H'{}^\dagger\) and \(I-\mathsf H^\dagger\mathsf H\) are orthogonal projection residuals, so their operator norms are at most one.  Taking norms term by term gives
\[
\begin{aligned}
\|\mathsf H^\dagger(\mathsf H-\mathsf H')\mathsf H'{}^\dagger\|_{\mathrm{op}}
&\le s_\star^{-2}\|\mathsf H-\mathsf H'\|_{\mathrm{op}},\\
\|\mathsf H^\dagger(\mathsf H^\dagger)^\top
  (\mathsf H-\mathsf H')^\top(I-\mathsf H'\mathsf H'{}^\dagger)\|_{\mathrm{op}}
&\le s_\star^{-2}\|\mathsf H-\mathsf H'\|_{\mathrm{op}},\\
\|(I-\mathsf H^\dagger\mathsf H)(\mathsf H-\mathsf H')^\top
  (\mathsf H'{}^\dagger)^\top\mathsf H'{}^\dagger\|_{\mathrm{op}}
&\le s_\star^{-2}\|\mathsf H-\mathsf H'\|_{\mathrm{op}}.
\end{aligned}
\]
Hence
\[
\|\mathsf H^\dagger-\mathsf H'{}^\dagger\|_{\mathrm{op}}
\le
3s_\star^{-2}\|\mathsf H-\mathsf H'\|_{\mathrm{op}} .
\]
For matrices \(\mathsf N,\mathsf N'\in\mathbb R^{d_z\times d_x}\), the algebraic split
\[
\mathsf H^\dagger\mathsf N-\mathsf H'{}^\dagger\mathsf N'
=
(\mathsf H^\dagger-\mathsf H'{}^\dagger)\mathsf N
+\mathsf H'{}^\dagger(\mathsf N-\mathsf N')
\]
therefore implies, whenever \(\|\mathsf N\|_{\mathrm{op}}\le L\),
\[
\|\mathsf H^\dagger\mathsf N-\mathsf H'{}^\dagger\mathsf N'\|_{\mathrm{op}}
\le
3s_\star^{-2}L\,\|\mathsf H-\mathsf H'\|_{\mathrm{op}}
+s_\star^{-1}\|\mathsf N-\mathsf N'\|_{\mathrm{op}} .
\]
Applying this estimate with
\[
\mathsf H=M_t(P),\quad \mathsf H'=M_t(Q),\quad
\mathsf N=N_t(P),\quad \mathsf N'=N_t(Q)
\]
gives
\[
\begin{aligned}
&\|M_t(P)^\dagger N_t(P)-M_t(Q)^\dagger N_t(Q)\|_{\mathrm{op}}\\
&\qquad\le
3s_\star^{-2}L\,\|M_t(P)-M_t(Q)\|_{\mathrm{op}}
+s_\star^{-1}\|N_t(P)-N_t(Q)\|_{\mathrm{op}} .
\end{aligned}
\]
Summing over \(t=0,1\) and using the definition of \(d_S\) yields
\[
\|\mathfrak D_P-\mathfrak D_Q\|_{\mathrm{op}}
\le
\{3s_\star^{-2}L+s_\star^{-1}\}\,
d_S(S(P),S(Q)).
\]
Moreover,
\[
\|m_X(P)-m_X(Q)\|_2\le d_S(S(P),S(Q)),
\qquad
\|m_X(Q)\|_2\le L,
\qquad
\|e_1\|_2=1 .
\]
% lean: ambientOperatorBridge_norm_ambientEffectOperator_sub_le

\item Let \(\varphi\) be the normalized attaining Kantorovich--Rubinstein potential for
\((\nu_P,\nu_Q)\) from \cref{obj:lem:kantorovich-rubinstein-attaining-potential}.  Then
\[
\varphi(0)=0,
\qquad
\operatorname{Lip}(\varphi)\le1,
\qquad
W_1(\nu_P,\nu_Q)
=
\left|
\int\varphi\,d\nu_P-\int\varphi\,d\nu_Q
\right|.
\]
Use the diagonalizations in \cref{eq:gap-free-model-diagonalization} for \(P\) and \(Q\), written
\[
\mathfrak D_P
=
\mathcal S_P\operatorname{diag}(\alpha_{P,1},\ldots,\alpha_{P,d_x})\mathcal S_P^{-1},
\qquad
\mathfrak D_Q
=
\mathcal S_Q\operatorname{diag}(\alpha_{Q,1},\ldots,\alpha_{Q,d_x})\mathcal S_Q^{-1}.
\]
Let
\[
\Omega_P:=\{\alpha_{P,j}:1\le j\le d_x\},
\qquad
\Omega_Q:=\{\alpha_{Q,j}:1\le j\le d_x\}.
\]
For any real number \(\zeta\), define the aggregate projectors
\[
\mathsf{Proj}_P(\zeta)
:=
\mathcal S_P\operatorname{diag}\{\mathbf 1(\alpha_{P,j}=\zeta):1\le j\le d_x\}\mathcal S_P^{-1},
\]
and
\[
\mathsf{Proj}_Q(\zeta)
:=
\mathcal S_Q\operatorname{diag}\{\mathbf 1(\alpha_{Q,j}=\zeta):1\le j\le d_x\}\mathcal S_Q^{-1}.
\]
The diagonal \(0\)-\(1\) matrices have operator norm at most one, so for
\(\zeta_P\in\Omega_P\) and \(\zeta_Q\in\Omega_Q\),
\[
\|\mathsf{Proj}_P(\zeta_P)\|_{\mathrm{op}}\le\kappa_\star,
\qquad
\|\mathsf{Proj}_Q(\zeta_Q)\|_{\mathrm{op}}\le\kappa_\star.
\]
The projectors extract their spectral values on the appropriate side:
\[
\mathsf{Proj}_P(\zeta_P)\mathfrak D_P
=
\zeta_P\mathsf{Proj}_P(\zeta_P),
\qquad
\mathfrak D_Q\mathsf{Proj}_Q(\zeta_Q)
=
\zeta_Q\mathsf{Proj}_Q(\zeta_Q).
\]
Hence
\[
\begin{aligned}
&\mathsf{Proj}_P(\zeta_P)(\mathfrak D_P-\mathfrak D_Q)\mathsf{Proj}_Q(\zeta_Q)\\
&\qquad=
(\zeta_P-\zeta_Q)\mathsf{Proj}_P(\zeta_P)\mathsf{Proj}_Q(\zeta_Q).
\end{aligned}
\]
In particular, when \(\zeta_P=\zeta_Q\),
\[
\mathsf{Proj}_P(\zeta_P)(\mathfrak D_P-\mathfrak D_Q)\mathsf{Proj}_Q(\zeta_Q)=0.
\]
Functional calculus with the two diagonalizations gives
\[
\varphi(\mathfrak D_P)
=
\sum_{\zeta_P\in\Omega_P}\varphi(\zeta_P)\mathsf{Proj}_P(\zeta_P),
\qquad
\varphi(\mathfrak D_Q)
=
\sum_{\zeta_Q\in\Omega_Q}\varphi(\zeta_Q)\mathsf{Proj}_Q(\zeta_Q).
\]
Since the aggregate projectors sum to the identity for each diagonalization, subtraction and expansion over
\(\Omega_P\times\Omega_Q\) yield
\[
\varphi(\mathfrak D_P)-\varphi(\mathfrak D_Q)
=
\sum_{\zeta_P\in\Omega_P}\sum_{\zeta_Q\in\Omega_Q}
\Xi(\zeta_P,\zeta_Q),
\]
where
\[
\Xi(\zeta_P,\zeta_Q)
:=
\begin{cases}
0,
&\zeta_P=\zeta_Q,\\[2mm]
\dfrac{\varphi(\zeta_P)-\varphi(\zeta_Q)}{\zeta_P-\zeta_Q}\,
\mathsf{Proj}_P(\zeta_P)(\mathfrak D_P-\mathfrak D_Q)\mathsf{Proj}_Q(\zeta_Q),
&\zeta_P\ne\zeta_Q .
\end{cases}
\]
The equal-spectral-value branch is zero by the preceding vanishing identity.  In the unequal branch, the one-Lipschitz property gives
\[
\left|
\frac{\varphi(\zeta_P)-\varphi(\zeta_Q)}{\zeta_P-\zeta_Q}
\right|\le1,
\]
and therefore
\[
\|\Xi(\zeta_P,\zeta_Q)\|_{\mathrm{op}}
\le
\kappa_\star^2\|\mathfrak D_P-\mathfrak D_Q\|_{\mathrm{op}}.
\]
Because \(|\Omega_P|\le d_x\) and \(|\Omega_Q|\le d_x\), the triangle inequality gives
\[
\|\varphi(\mathfrak D_P)-\varphi(\mathfrak D_Q)\|_{\mathrm{op}}
\le
d_x^2\kappa_\star^2
\|\mathfrak D_P-\mathfrak D_Q\|_{\mathrm{op}}.
\]
Also,
\[
\|\varphi(\mathfrak D_P)\|_{\mathrm{op}}
\le \kappa_\star L_\tau,
\]
because \(\varphi(0)=0\), \(\varphi\) is one-Lipschitz, and every eigenvalue of \(\mathfrak D_P\) lies in \([-L_\tau,L_\tau]\).

Using the representation identity for \(P\) and \(Q\),
\[
\begin{aligned}
W_1(\nu_P,\nu_Q)
&\le
\|m_X(P)-m_X(Q)\|_2\,\|\varphi(\mathfrak D_P)\|_{\mathrm{op}}\|e_1\|_2\\
&\quad+
\|m_X(Q)\|_2\,
\|\varphi(\mathfrak D_P)-\varphi(\mathfrak D_Q)\|_{\mathrm{op}}\|e_1\|_2\\
&\le
C_{\mathrm{base}}\,d_S(S(P),S(Q))
\le
C_{\mathrm{mod}}\,d_S(S(P),S(Q)).
\end{aligned}
\]
This proves the asserted modulus on model summaries:
\begin{equation}
W_1(\nu_P,\nu_Q)
\le C_{\mathrm{mod}}d_S(S(P),S(Q)).
\label{eq:gap-free-model-summary-control}
\end{equation}
% lean: ambientOperatorBridge_modelLaw_wass1_le_dS_of_certificates

\item Let
\[
\mathscr S:=\{S(P):P\in\mathcal M\}.
\]
For \(q\in\mathscr S\), choose a representative law \(P_q\in\mathcal M\) with
\[
S(P_q)=q,
\]
and define
\[
F_{\mathscr S}(q):=\nu_{P_q}.
\]
This definition is independent of the representative.  Indeed, if another \(Q\in\mathcal M\) has \(S(Q)=q\), then the modulus just proved gives
\[
0\le W_1(\nu_{P_q},\nu_Q)
\le C_{\mathrm{mod}}\,d_S(q,q)=0.
\]
A zero finite transport cost between quotient atomic laws forces the same aggregated mass at each support point, so the two quotient laws are equal.  Hence
\[
F_{\mathscr S}(S(Q))=\nu_Q
\qquad(Q\in\mathcal M),
\]
and, for all \(q,q'\in\mathscr S\),
\[
W_1(F_{\mathscr S}(q),F_{\mathscr S}(q'))
\le
C_{\mathrm{mod}}\,d_S(q,q').
\]
% lean: gap_free_positive_measure_modulus

\item Let \(\rho_{\mathrm{sum}}\) be the ambient product metric on the five summary blocks: each matrix block has its Frobenius norm, the mean block has its Euclidean norm, and the finite product uses the maximum of the five block distances.  Set
\[
\Lambda_{\mathrm{sum}}:=4\operatorname{Ent}_{d_z,d_x}+d_x .
\]
By \cref{obj:lem:summary-metric-coordinate-comparison},
\[
d_S(q,q')
\le
\Lambda_{\mathrm{sum}}\rho_{\mathrm{sum}}(q,q')
\qquad\text{for all summaries }q,q'.
\]
For a matrix block \(A\in\mathbb R^{d_z\times d_x}\), each entry obeys
\[
|A_{ab}|\le\|A\|_{\mathrm{op}},
\]
so summing the squared entries gives
\[
\|A\|_{\mathrm F}\le\sqrt{d_zd_x}\,\|A\|_{\mathrm{op}}.
\]
Since \(\rho_{\mathrm{sum}}\) is the maximum of the four Frobenius block distances and the Euclidean mean-block distance, while \(d_S\) is the sum of the four operator-norm distances and that same mean-block distance,
\[
\rho_{\mathrm{sum}}(q,q')
\le
\max\{1,\sqrt{d_zd_x}\}\,d_S(q,q')
\qquad\text{for all summaries }q,q'.
\]
(The assumptions \(2\le k\le d_x,d_z\) cover the positive-dimensional case used here; in a zero-dimensional matrix block the unique empty matrix contributes zero.)  Thus \(d_S\) and \(\rho_{\mathrm{sum}}\) induce the same topology and the same closures.

The first comparison and the \(d_S\)-control on \(F_{\mathscr S}\) give
\[
W_1(F_{\mathscr S}(q),F_{\mathscr S}(q'))
\le
C_{\mathrm{mod}}\Lambda_{\mathrm{sum}}\rho_{\mathrm{sum}}(q,q')
\qquad(q,q'\in\mathscr S).
\]
Hence \(F_{\mathscr S}\) is Lipschitz on the summary image in the ambient metric.

Inside \(\mathcal K=\overline{\mathscr S}\), put
\[
\mathscr S_{\mathcal K}:=\{q\in\mathcal K:q\in\mathscr S\}.
\]
The equality of closures just proved makes \(\mathscr S_{\mathcal K}\) dense in \(\mathcal K\).  The quotient atomic-law space \(\mathcal P_k([-L_\tau,L_\tau])\) is complete for \(W_1\): in labelled coordinates it is the metric quotient of the compact set
\[
\Delta_{k-1}\times[-L_\tau,L_\tau]^k.
\]
Apply the dense-set extension theorem to the ambient-Lipschitz map
\[
\widetilde F_{\mathscr S}(q):=F_{\mathscr S}(q),
\qquad q\in\mathscr S_{\mathcal K}.
\]
The dense-set extension construction gives a Lipschitz, hence continuous, map
\[
\overline F:\mathcal K\to\mathcal P_k([-L_\tau,L_\tau])
\]
that agrees with \(F_{\mathscr S}\) on \(\mathscr S\), and the uniqueness property of this construction makes it the unique continuous extension with that agreement.

It remains to transfer the sharper \(d_S\)-control.  Fix \(q,q'\in\mathcal K\).  Choose sequences \(q_m,q'_m\in\mathscr S\) converging in the ambient metric to \(q,q'\), respectively.  Continuity of \(\overline F\) and \cref{obj:lem:summary-metric-continuity} give
\[
W_1(\overline F(q_m),\overline F(q'_m))
\longrightarrow
W_1(\overline F(q),\overline F(q')),
\qquad
d_S(q_m,q'_m)\longrightarrow d_S(q,q').
\]
For every \(m\), agreement on \(\mathscr S\) and the model-summary control in \cref{eq:gap-free-model-summary-control} yield
\[
W_1(\overline F(q_m),\overline F(q'_m))
=
W_1(F_{\mathscr S}(q_m),F_{\mathscr S}(q'_m))
\le
C_{\mathrm{mod}}d_S(q_m,q'_m).
\]
Passing to the limit proves
\[
W_1(\overline F(q),\overline F(q'))
\le
C_{\mathrm{mod}}d_S(q,q')
\qquad(q,q'\in\mathcal K).
\]
Continuity of \(\overline F\), together with the Borel structure on the quotient atomic-law space, makes \(\overline F\) Borel measurable.
% lean: gapFreeClosureAssembly_exists_unique_extension_with_control

\item For every \(Q\in\mathcal M\), the summary \(S(Q)\) belongs to \(\mathscr S\), and hence
\[
\overline F(S(Q))
=
F_{\mathscr S}(S(Q))
=
\nu_Q .
\]
Let
\[
G:\mathcal K\to\mathcal P_k([-L_\tau,L_\tau])
\]
be continuous, satisfy the stated \(C_{\mathrm{mod}}\)-Lipschitz bound, and agree with \(Q\mapsto\nu_Q\) on all model summaries.  For any \(q\in\mathcal K\), choose \(q_m\in\mathscr S\) with
\[
q_m\to q .
\]
Continuity and agreement on \(\mathscr S\) imply
\[
G(q)
=
\lim_m G(q_m)
=
\lim_m F_{\mathscr S}(q_m)
=
\overline F(q).
\]
Thus \(G(q)=\overline F(q)\) for all \(q\in\mathcal K\).  Combining the model-summary modulus, the continuous extension, the Borel measurability, and this uniqueness proves the theorem.
% lean: gap_free_positive_measure_modulus
\end{enumerate}
\end{proof}
